The parietal and occipital electrodes performed markedly worse than the frontal and central electrodes, and their performance differed between datasets. On Internal, the four parietal electrodes, CP1, CP2, CP5 and CP6, were uniformly near floor, with $F_1$ values between 1.75\% and 2.19\%. The eight occipital electrodes, P7, P3, PO3, O1, PO4, E83, P4 and P8, showed slightly more spread, from 1.77\% at P4 to 3.60\% at PO3. On Cao2018, the two parietal electrodes were more informative but remained well below the frontal and central regions, with CP3 and CP4 reaching 17.84\% and 19.56\%. The Cao2018 occipital region was also more heterogeneous than its Internal counterpart. P3 and P4 scored highest at 12.80\% and 12.67\%, T5 and T6 were intermediate at 9.26\% and 8.30\%, and O1 and O2 were near floor at 3.10\% and 3.33\%.

\begin{figure}[ht]
    \centering
    \includegraphics[width=0.9\linewidth]{fig_exp1_full_montage_per_electrode.pdf}
    \caption{Per-electrode detection $F_1$ for every electrode in the four curated
        anatomical regions, scored in the full-montage run, coloured by scalp region. Each bar is one electrode's session-mean $F_1$
        under the complete pipeline using all 32 channels.
%        ; this is the full-montage
%        baseline that the regional
%        (Figure~\ref{fig:exp1_subset_per_electrode})
%        and single-channel
%        (Figure~\ref{fig:exp1_single_channel})
%        reruns are compared
%        against.
    }
    \label{fig:exp1_full_montage_per_electrode}
\end{figure}
